% Hafta 6 — WAF ile RASP farkı
% Derleme: py -3.12 tools/latex/derle.py docs/week-6/assets/h06-06-waf-rasp.tex
\documentclass[tikz,border=8pt]{standalone}
\usepackage{cen429-cizim}
\begin{document}
\begin{tikzpicture}
  \node[baslik] (b) at (0,0) {WAF ile RASP farkı};
  \node[altbaslik] at (b.south west) {``kurcalandım mı, beni kim izliyor, hangi cihazdayım?''};
  \node[kutu=70mm, anchor=north west] (sol) at (0,-2.0)
    {\kb{WAF}\\[2mm]
     uygulamanın ÖNÜNDE\\HTTP isteklerini süzer\\[3mm]
     uygulamanın içini bilmez\\sunucu tarafı};
  \node[iyikutu=70mm, right=14mm of sol] (sag)
    {\kb{RASP}\\[2mm]
     uygulamanın İÇİNDE\\kendi kodunu, belleğini,\\ortamını izler\\[3mm]
     MATE saldırganına cevap};
  \node[fit=(sol)(sag), inner sep=0pt] (grp) {};
  \node[dip, text width=155mm, anchor=north west] at ($(grp.south west)+(0,-8mm)$)
    {RASP, cihazın güvenilmez olduğu durumlar için tasarlanmıştır.};
\end{tikzpicture}
\end{document}
